\section{Scriptless Escrow}
\label{sec:escrow}

A moderated purchase protects the buyer against non-delivery and the vendor against a buyer who will not release payment, with a moderator to decide disputes. The original OpenBazaar used a 2-of-3 multisignature script for this. The redesign keeps the economic model (a moderator together with one party can decide the outcome) and changes three things: the moderator holds no key over the funds; every outcome looks like an ordinary single-key Taproot spend; and the moderator can choose only between two payouts that buyer and vendor both signed before funding.

This section first states what is merged and what the open escrow pull request changes, then specifies the two-party protocol.

\subsection{Status of the escrow code}
\label{sec:escrow-status}

\paragraph{Merged (main).} The escrow crate on main computes a MuSig2 key-path funding output, presigns a locktime refund, and builds adaptor signatures for both dispute outcomes \ID. It does all of this in one function, \code{fund}, that takes both the buyer's and the vendor's secret keys. A review of the merged code found that, in addition: the node calls it with fixture keys; the funding outpoint is synthetic; the presigned refund is not kept by the node; there is no presigned vendor default, so a buyer who stops responding cannot be overridden before the refund height; the encrypted moderator halves carry no proof that they decrypt to the right value; there is no moderator fee; network fees are zero; and the funding address is the same for every order between the same buyer and vendor keys. The merged escrow is a demonstration of the transaction shapes, not a protocol between two parties.

\paragraph{PR~\#2.} The payments branch keeps the single-process ceremony but restricts it to testnet4 (mainnet refuses escrow), derives a new escrow key per order for each side, adds a 500-satoshi spend fee, and enforces the order ``sign the funding transaction, run the ceremony for that exact outpoint, check and store the refund, then broadcast'' \ID. Its responses mark escrow as an unsafe demonstration.

\paragraph{PR~\#3.} The two-party protocol below is implemented \IT{} in the escrow crate on branch \code{escrow-two-party} (we ran its 26 unit tests and 3 doctests). Each party's state machine holds only that party's secret key, and the parties exchange three serializable messages. The node does not call this protocol yet, no transport carries its messages, and the old single-process function remains behind a deprecated feature so the node demonstration still builds. Everything in \S\ref{sec:escrow-notation}--\S\ref{sec:escrow-privacy} describes PR~\#3; we note where its behavior depends on a caller obligation that no code yet fulfills.

\subsection{Notation}
\label{sec:escrow-notation}

$G$ generates the secp256k1 group of prime order $n$. Lowercase letters are scalars and uppercase letters are points. $\Hash_{\mathrm{tag}}$ is a BIP 340 tagged hash.

\begin{itemize}
  \item Long-term keys: buyer $(\bar b,\bar B)$, vendor $(\bar v,\bar V)$. Moderator $(k,K)$.
  \item Order secret $s$: 32 bytes known only to buyer and vendor. Per-order keys $b=\bar b+\Hash_{\mathrm{ord}}(\bar B,s)$ and $B=bG$, and likewise $v$, $V$.
  \item Aggregate key: $P=\mathsf{KeyAgg}(B,V)$ under BIP 327~\citep{bip327} with the two keys sorted; output key $Q=P+\Hash_{\mathrm{TapTweak}}(P)\,G$, the BIP 341~\citep{bip341} tweak with no script tree. The funding output is P2TR with output key $Q$ and value $F$.
  \item Terms $\tau=(B,V,K,F,f,\varphi,h_D,h_C,h_R,S_B,S_V,S_M)$: network fee $f$, moderator fee $\varphi$, dispute deadline $h_D$, vendor claim height $h_C$, refund height $h_R$, and payout scripts for buyer, vendor, and moderator. Both parties sign $\tau$; the signed terms are $\tau^\ast$.
  \item Outpoint $o$ of the funding output.
  \item For a dispute branch favoring party $X\in\{B,V\}$, with $\bar X$ the other party: the \emph{party half} $t_X$, $T_X=t_XG$, derived by $X$ from its per-order secret key and the branch; the \emph{moderator half} $u_X$, $U_X=u_XG$, sampled by $\bar X$; the adaptor point $Y_X=T_X+U_X$.
  \item $\mathsf{VEnc}_K(u;\mathit{ctx})$: verifiable encryption of $u$ to $K$, bound to context $\mathit{ctx}$ (\S\ref{sec:vencrypt}). $\pi^{\mathrm{pok}}(t)$: Schnorr proof of knowledge of $t$.
\end{itemize}

\begin{figure}[t]
\centering
\begin{tikzpicture}[
  font=\small,
  txb/.style={draw, rounded corners=2pt, align=left, text width=5.6cm, minimum height=1.0cm, fill=white, inner sep=4pt},
  fund/.style={draw, rounded corners=2pt, align=center, text width=3.6cm, minimum height=1.6cm, fill=black!5, font=\footnotesize},
  arr/.style={-{Stealth[length=2mm]}, thick},
]
\node[fund] (f) at (0,0) {\textbf{Funding output}\\[2pt] P2TR, key path only\\ $Q = P + \Hash_{\mathrm{TapTweak}}(P)\,G$\\ $P=\mathsf{KeyAgg}(B,V)$\\ value $F$};
\node[txb] (rel) at (6.8,3.6)  {\textbf{Release} \hfill $\mathsf{lt}=0$\\ $F-f \to$ vendor\\[-1pt]\scriptsize signed after delivery};
\node[txb] (clm) at (6.8,1.8) {\textbf{Vendor claim} \hfill $\mathsf{lt}=h_C$\\ $F-f-a \to$ vendor, anchor\\[-1pt]\scriptsize buyer partial in round 3};
\node[txb] (ref) at (6.8,0)    {\textbf{Buyer refund} \hfill $\mathsf{lt}=h_R$\\ $F-f-a \to$ buyer, anchor\\[-1pt]\scriptsize vendor partial in round 2};
\node[txb] (bw)  at (6.8,-1.8){\textbf{Buyer wins} \hfill $\mathsf{lt}=0$\\ $F-f-a-\varphi \to$ buyer, $\varphi\to$ mod., anchor\\[-1pt]\scriptsize adaptor $T_B+U_B$, round 2};
\node[txb] (vw)  at (6.8,-3.6) {\textbf{Vendor wins} \hfill $\mathsf{lt}=0$\\ $F-f-a-\varphi \to$ vendor, $\varphi\to$ mod., anchor\\[-1pt]\scriptsize adaptor $T_V+U_V$, round 3};
\foreach \t in {rel,clm,ref,bw,vw} { \draw[arr] (f.east) -- (\t.west); }
\node[font=\scriptsize, align=center, text width=3.4cm] at (0,-2.0) {witness of every spend:\\one 64-byte BIP~340 signature};
\end{tikzpicture}
\caption{Escrow transaction graph. All five spends are key-path spends of the same output under the MuSig2 aggregate of the buyer's and vendor's per-order keys; the moderator's key is not part of $Q$. Four spends are presigned before funding; the release is signed cooperatively later. $f$ is the network fee, $\varphi$ the moderator fee, and $a$ the 240-satoshi pay-to-anchor output, all fixed by the signed terms; $\mathsf{lt}$ is \code{nLockTime}.}
\label{fig:escrow-graph}
\end{figure}

\begin{figure}[t]
\centering
\begin{tikzpicture}[font=\small, >={Stealth[length=2mm]}]
\draw[->, thick] (0,0) -- (14,0) node[right]{height};
\foreach \x/\l in {0.4/funding, 5.5/$h_D$, 9.0/$h_C$, 12.5/$h_R$} {
  \draw[thick] (\x,0.15) -- (\x,-0.15);
  \node[below] at (\x,-0.15) {\l};
}
\draw[|-|] (0.4,1.25) -- (5.5,1.25);
\node[above, align=center, font=\scriptsize] at (2.95,0.4) {ship, deliver ($\ge$1008 blocks);\\dispute letters accepted};
\node[above, align=center, font=\scriptsize] at (7.25,0.4) {ruling window\\($\ge$144 blocks)};
\node[above, align=center, font=\scriptsize] at (10.75,0.4) {claim window\\($\ge$144 blocks)};
\draw[|-|] (5.5,1.25) -- (9.0,1.25);
\draw[|-|] (9.0,1.25) -- (12.5,1.25);
\node[below, align=center, font=\scriptsize] at (9.0,-0.65) {vendor claim\\minable at $h_C+1$};
\node[below, align=center, font=\scriptsize] at (12.5,-0.65) {buyer refund\\minable at $h_R+1$};
\draw[dashed] (0.4,-1.65) -- (13.6,-1.65);
\node[font=\scriptsize, align=left, anchor=west] at (0.4,-2.0) {release and both dispute outcomes are valid at every height once their secrets are known};
\end{tikzpicture}
\caption{Escrow timeline. The default after $h_C$ is the vendor; the refund at $h_R$ matters only if the vendor never claims. Window sizes are policy minimums checked by each party at funding against the current tip; the minimums are 1{,}008 blocks from the tip to $h_D$ and 144 blocks for each later window.}
\label{fig:escrow-timeline}
\end{figure}


\subsection{Terms and transactions}

Before any funding round, the parties agree on $\tau$ and each checks it under its own local policy: the three keys are distinct; $0<h_D<h_C<h_R$ are block heights with $h_C-h_D$ and $h_R-h_C$ at least a minimum window (144 blocks each by default); the moderator fee is at least the P2TR dust limit of 330 satoshis; every payout is at least dust; and every presigned transaction meets a minimum fee rate measured on its actual signed size \IT. The moderator key must match the key the party already trusts from the listing, so a counterparty cannot substitute a different moderator \IT.

From $\tau$ and $o$, each party derives the same five spends (Figure~\ref{fig:escrow-graph}). Each has one input spending $o$ with sequence \code{0xfffffffe} (so \code{nLockTime} is enforced) and is signed as a BIP 341 key-path spend with \code{SIGHASH\_DEFAULT}:
\begin{center}\small
\begin{tabular}{@{}llll@{}}
\toprule
\textbf{Spend} & \code{nLockTime} & \textbf{Outputs} & \textbf{When signed} \\
\midrule
$\tx{rel}$ (release) & 0 & $F-f\to S_V$ & cooperatively, after delivery \\
$\tx{clm}$ (vendor claim) & $h_C$ & $F-f\to S_V$ & presigned (buyer, round 3) \\
$\tx{ref}$ (buyer refund) & $h_R$ & $F-f\to S_B$ & presigned (vendor, round 2) \\
$\tx{B}$ (buyer wins) & 0 & $F-f-\varphi\to S_B$, $\varphi\to S_M$ & adaptor $Y_B$ (vendor, round 2) \\
$\tx{V}$ (vendor wins) & 0 & $F-f-\varphi\to S_V$, $\varphi\to S_M$ & adaptor $Y_V$ (buyer, round 3) \\
\bottomrule
\end{tabular}
\end{center}
A transaction with \code{nLockTime} $h$ can first be mined in block $h+1$.

\subsection{Funding rounds}

Let $\mathsf{N}$ denote a party's MuSig2 public nonces for the four presigned spends and $\mathsf{ps}(\cdot)$ a MuSig2 partial signature~\citep{nick2021musig2} (or partial adaptor signature, for $\tx{B}$ and $\tx{V}$).

\begin{description}[leftmargin=1.2em]
  \item[Round 0.] Both parties sign $\tau$.
  \item[Round 1, buyer $\to$ vendor.] The buyer prepares, but does not sign or broadcast, a funding transaction from native segwit inputs, so its transaction id is fixed before signing. It sends
  $\bigl(o,\ \mathsf{N}_B,\ T_B,\ \pi^{\mathrm{pok}}(t_B),\ \mathsf{VEnc}_K(u_V;\mathit{ctx}_V)\bigr)$.
  The buyer samples $u_V$ because the buyer loses branch $V$.
  \item[Round 2, vendor $\to$ buyer.] The vendor checks $\pi^{\mathrm{pok}}(t_B)$ and the encryption proof for $U_V$ (it is the favored party in branch $V$, so it must be sure the moderator can reveal $u_V$). It replies
  $\bigl(\mathsf{N}_V,\ T_V,\ \pi^{\mathrm{pok}}(t_V),\ \mathsf{VEnc}_K(u_B;\mathit{ctx}_B),\ \mathsf{ps}_V(\tx{ref}),\ \mathsf{ps}_V(\tx{B};Y_B)\bigr)$.
  \item[Round 3, buyer $\to$ vendor.] The buyer checks $\pi^{\mathrm{pok}}(t_V)$ and the encryption proof for $U_B$, adds its own partial signatures, and verifies that the completed refund is a valid signature under $Q$ for the exact outpoint and value, and that the buyer-wins pre-signature is valid for $Y_B$. Only after these checks does the buyer's state machine release the funding transaction for signing \IT. The buyer must persist the refund before funding; the crate exposes it for storage and re-verification, and storing it is the caller's obligation. The buyer sends
  $\bigl(\mathsf{ps}_B(\tx{clm}),\ \mathsf{ps}_B(\tx{V};Y_V)\bigr)$.
  The vendor completes and stores the claim and verifies the vendor-wins pre-signature.
  \item[Funding.] The buyer signs and broadcasts the funding transaction. The vendor's state machine checks that a transaction creates exactly the agreed output at $o$; shipping only after that transaction confirms is the caller's obligation.
\end{description}

Each state is consumed by value when it advances, and secret nonces live only inside those states, so a nonce cannot be reused for two messages by accident. The release uses fresh nonces generated later.

\paragraph{Why each piece is there.}
\begin{itemize}
  \item \emph{Refund before funding.} The buyer never signs the funding transaction until it holds a refund that verifies against the aggregate key. A vendor who aborts after round 1 costs the buyer nothing.
  \item \emph{Claim before refund.} The buyer presigns a claim at $h_C<h_R$, so a buyer who received the goods and stops responding loses at $h_C+1$.
  \item \emph{Proof of knowledge on $T_X$.} Without it, $X$ could choose $T_X=T'-U_X$ for a $T'$ whose discrete log it knows, making $Y_X=T'$ and letting $X$ complete its own winning branch without the moderator.
  \item \emph{Verifiable encryption of $u_X$.} Without it, the losing party $\bar X$ could encrypt a value that does not match $U_X$, and the moderator would be unable to rule for $X$.
  \item \emph{Per-order keys.} Tweaking each long-term key with the order secret gives a different $Q$ for every order between the same two parties, so funding outputs between repeat trading partners are not linkable by address \IT. The protocol assumes $s$ is shared only by buyer and vendor; carrying it inside the sealed order letter is the intended channel \SP.
\end{itemize}

\subsection{Release}

When the buyer is satisfied, the vendor sends a fresh public nonce for $\tx{rel}$, and the buyer returns its nonce and partial signature. The vendor aggregates them into a 64-byte signature and broadcasts \IT. The moderator is not involved.

\subsection{Disputes}

Before $h_D$, either party $X$ can ask the moderator to rule for it by sending a \emph{dispute letter} \IT:
\[
  L_X = \bigl(\tau^\ast,\ o,\ F,\ \mathit{script}(Q),\ Q,\ X,\ \tx{X}\ \text{(unsigned)},\ T_X,\ \mathsf{VEnc}_K(u_X),\ \hat\sigma_X,\ \mathit{memo}\bigr)\ \text{signed by } X.
\]
Here $\hat\sigma_X$ is the aggregated adaptor pre-signature on $\tx{X}$ for $Y_X$. The letter is self-verifying: the moderator checks the signed terms against its own key, recomputes $P$, $Q$, the funding script, and $\tx{X}$ from the terms and outpoint, checks that $\hat\sigma_X$ is a valid pre-signature on $\tx{X}$ for $T_X+U_X$, checks the encryption proof, and checks $X$'s signature over the letter. It refuses a letter that arrives after $h_D$. Only then does it decrypt $u_X$.

If it rules for $X$, the moderator returns a \emph{ruling} $(o, X, Y_X, u_X)$ signed with $k$. $X$ computes $t_X+u_X$, completes $\hat\sigma_X$ into a valid BIP 340 signature, and broadcasts $\tx{X}$, which pays $X$ and the moderator's fee.

A ruling for $X$ is useless to $\bar X$, because $\bar X$ does not know $t_X$. A moderator that signs rulings for both branches of the same outpoint produces two signed rulings that together are transferable proof of equivocation \IT. Equivocation is detectable, not preventable: both dispute outcomes are valid from height 0, so if both parties receive a ruling, the first to confirm wins.

\subsection{Timeline}

Figure~\ref{fig:escrow-timeline} shows the heights.
\begin{itemize}
  \item From funding until $h_D$, the vendor ships, the buyer receives, and either party may send a dispute letter.
  \item Between $h_D$ and $h_C$, the moderator rules and the winner must confirm its outcome. A buyer who wins must confirm $\tx{B}$ before $h_C+1$, when $\tx{clm}$ also becomes valid.
  \item At $h_C+1$, the vendor can claim. A vendor who never claims leaves the funds in the output until $h_R+1$, when the buyer's refund becomes valid. If the vendor's claim is not confirmed by then, both are valid and the first to confirm wins.
\end{itemize}

The default outcome is therefore the vendor's. We considered and rejected a design in which the default flips to the buyer until the vendor posts a shipping attestation: without a trusted third party that observes shipping, such an attestation is not a cryptographic event, and a buyer who wants to cheat would decline to countersign it. The consequence is that a buyer who did not receive the goods must open a dispute before $h_D$.

\subsection{Moderator fee}
The moderator fee is an output of each dispute outcome and of nothing else, and it is fixed in the signed terms \IT. It is collected only when the winner broadcasts the dispute outcome. A vendor who wins a ruling close to $h_C$ can skip the fee by waiting for the fee-free claim; a buyer who wins cannot. That skews moderator income toward rulings for the buyer. A fee bond posted at funding would remove the skew and is not implemented \FW.

\subsection{Verifiable encryption of the moderator half}
\label{sec:vencrypt}

The moderator half $u$ with $U=uG$ is encrypted bit by bit with exponential ElGamal under $K$ \IT. For each bit $\beta_i$ of $u$ ($i=0,\dots,255$), the encryptor picks $r_i$ and publishes $(A_i,C_i)=(r_iG,\ r_iK+\beta_iG)$ with a Cramer--Damg{\aa}rd--Schoenmakers OR proof~\citep{cramer1994} of two Chaum--Pedersen statements~\citep{chaum1992} showing that $\beta_i\in\{0,1\}$. A Chaum--Pedersen proof then shows
\[
  \log_G\Bigl(\textstyle\sum_i 2^iA_i\Bigr)=\log_K\Bigl(\textstyle\sum_i 2^iC_i-U\Bigr),
\]
which holds exactly when $\sum_i 2^i\beta_i\equiv u \pmod n$. Anyone with $K$ and $U$ can check the ciphertext; only the holder of $k$ can decrypt, since $C_i-kA_i$ is either the identity or $G$. All challenges are tagged hashes over length-prefixed inputs and a 32-byte context that binds the proof to one escrow branch. This is the bitwise approach to verifiable encryption of a discrete logarithm~\citep{stadler1996pvss,camenisch2003}. Its cost is size: one encrypted share is 49{,}761 bytes (\S\ref{sec:evaluation}).

\subsection{What the chain shows}
\label{sec:escrow-privacy}

Every spend is a single-input Taproot key-path spend with one 64-byte signature. No script, no moderator key, and no multisignature structure appears on chain in any outcome, including a dispute \IT. This is the main privacy improvement over the original 2-of-3 design. What remains visible:
\begin{itemize}
  \item The funding amount $F$ (price plus shipping plus fees), which can point to a listing with a distinctive price.
  \item Dispute outcomes have two outputs, one of which is the fixed moderator fee $\varphi$. If a moderator reuses its payout script $S_M$, every dispute it rules on is linked on chain to it and to the others. Fresh moderator payout scripts per order are a caller obligation, not enforced.
  \item The claim and refund carry non-zero \code{nLockTime} values at specific heights, and all spends signal no replace-by-fee. These are weak fingerprints, since many wallets set \code{nLockTime} near the current height.
  \item Timing: a spend shortly after funding versus one after $h_C$.
\end{itemize}

\subsection{Operational limits}

Presigned transactions have fees fixed at funding. If fee rates rise sharply, a presigned spend may not confirm within its window. The winner of any outcome controls an output and could bump the fee by spending that output (child-pays-for-parent); no wallet code does this yet. Each verifiable share is about 50~KB, so a dispute letter is close to the 64~KiB limit of the order mailbox; the transport for escrow messages and dispute letters is not specified in detail. The protocol envelope used for disputes (\code{DisputeEnvelope} in \code{crates/protocol}) currently embeds the full order plaintext, including the ship-to field, so a moderator who receives it would learn the address; stripping the ship-to from dispute envelopes is required by the design \SP{} and not enforced.
